\begin{helper}
Fix a history event \(H\), a target event \(T\), a terminal coordinate set
\(U\), a nodup terminal order \(o\) with \(\operatorname{toFinset}(o)=U\),
a finite blue-label family \({\mathcal B}\), branch labels
\(\beta\), red labels \(J\), transcripts
\(\tau=(P_\tau,A_\tau,Z_\tau)\), blue traces \(S_\beta\), red supports
\(R_{\beta,J}\), and cell events \(G_{B,J,\beta,\tau}\).  Assume
\[
0\le p\le 1/2,
\]
and assume the complete terminal product law conditioned on \(H\): for every
assignment \(A\subseteq U\),
\[
\Pr(\forall e\in U,\ X_e\leftrightarrow e\in A\mid H)
=p^{|A|}(1-p)^{|U|-|A|}.
\]
Assume each \(B\in{\mathcal B}\) is a blue \(u_B\)-set contained in \(U\),
and assume branch functionality: equality of the complete blue trace on all
blue coordinates of \(U\) forces equality of branch labels.

For every nonempty cell, assume the branch/transcript geometry supplied by
\(\noderef{FixedSetHistoryCellRedBranchTranscriptCells}\):
\[
|R_{\beta,J}|=u_R,\qquad R_{\beta,J}\subseteq U,
\]
\[
R_{\beta,J}\hbox{ is red},\qquad
P_\tau,A_\tau\subseteq U,\qquad P_\tau\cap A_\tau=\emptyset,
\]
\[
B\cup R_{\beta,J}\subseteq P_\tau,\qquad
Z_\tau\subseteq A_\tau,\qquad
L:=\max(0,a-u_R-u_B)\le |Z_\tau|,
\]
and \(P_\tau,A_\tau\) encode the complete blue terminal trace: for every
blue coordinate \(e\in U\),
\[
e\in S_\beta\Longleftrightarrow e\in P_\tau,\qquad
e\notin S_\beta\Longleftrightarrow e\in A_\tau.
\]
the sample in the cell agrees with \(S_\beta\) on all blue terminal
coordinates, all coordinates of \(P_\tau\) are present, and all coordinates
of \(A_\tau\) are absent.  The selected absence set \(Z_\tau\) may contain
both red and blue coordinates.

In addition, assume the paper-specific combinatorial support interface
exposed by the branch/transcript construction.  There is a realized-cell
predicate \(\operatorname{Realized}(B,J,\beta,\tau)\) and an
assignment-compatibility predicate
\(\operatorname{Compat}(A;B,J,\beta,\tau)\).  The interface also includes an
explicit realized-geometry certificate: whenever
\(\operatorname{Realized}(B,J,\beta,\tau)\) holds, all the displayed
red-support, transcript-subset, disjointness, blue-trace, selected-absence,
and lower-bound properties above hold for that formal pair.  The
compatibility predicate implies the usual span equations: compatible
assignments contain \(P_\tau\), avoid \(A_\tau\setminus Z_\tau\), and agree
with \(S_\beta\) on every unselected blue coordinate.  Functional scan
disjointness says that one terminal assignment cannot be compatible with two
distinct realized pairs \((\beta,\tau)\).  Non-realized formal pairs have no
cell samples.
No one-cell released-reconstruction clause is assumed.  Instead, for every
fixed \((B,J)\), assume the residual leaf certificate supplied by the
lower-level reveal-order construction.  Thus there is a finite leaf set
\(\Lambda_{B,J}\), a leaf assignment
\[
\iota_{B,J}(\beta,\tau)\in\Lambda_{B,J},
\]
and a leaf mass \(\nu_{B,J}:\Lambda_{B,J}\to\mathbb R\) such that
\[
0\le \nu_{B,J}(\lambda),\qquad
\sum_{\lambda\in\Lambda_{B,J}}\nu_{B,J}(\lambda)\le1.
\]
The leaf assignment is injective on realized cells for this fixed pair:
if \(\tau,\tau'\) are allowed transcripts, both
\((B,J,\beta,\tau)\) and \((B,J,\beta',\tau')\) are realized, and
\[
\iota_{B,J}(\beta,\tau)=\iota_{B,J}(\beta',\tau'),
\]
then \(\beta=\beta'\) and \(\tau=\tau'\).  Finally, every realized pair
satisfies the residual domination
\[
p^{|P_\tau|-u_R-u_B}(1-p)^{|A_\tau|-|Z_\tau|}
\le \nu_{B,J}(\iota_{B,J}(\beta,\tau)).
\]
These are leaf-level scan data, not cell-indexed residual masses.

Then there are weights \(w_{B,J,\beta,\tau}\) such that
\[
0\le w_{B,J,\beta,\tau},
\qquad
w_{B,J,\beta,\tau}\le p^{|P_\tau|}(1-p)^{|A_\tau|},
\]
\[
\Pr(T\cap H\cap G_{B,J,\beta,\tau})
\le \Pr(H)\,w_{B,J,\beta,\tau},
\]
and, for every fixed \(B\in{\mathcal B}\) and \(J\),
\[
\sum_{\beta,\tau} w_{B,J,\beta,\tau}
\le p^{u_R}p^{u_B}(1-p)^L .
\]
\end{helper}
\begin{proof}
Write \(q=1-p\).  From \(0\le p\le1/2\), we have \(0\le q\le1\).  For a
formal pair \((B,J,\beta,\tau)\), define \(w_{B,J,\beta,\tau}=0\) if
\(\tau\) is not an allowed transcript or if
\(\operatorname{Realized}(B,J,\beta,\tau)\) fails.  For realized allowed
pairs, define \(w_{B,J,\beta,\tau}\) to be the terminal cylinder mass of the
cell block.

Nonnegativity follows immediately from the nonnegativity of the product-law
summands.  For a realized cell, the geometry gives
\[
P_\tau,A_\tau\subseteq U,\qquad P_\tau\cap A_\tau=\emptyset .
\]
Applying \(\noderef{FixedSetHistoryCellFixedCylinderCharge}\) to the present
set \(P_\tau\), the absent set \(A_\tau\), and the complete product law on
\(U\) gives
\[
w_{B,J,\beta,\tau}\le p^{|P_\tau|}q^{|A_\tau|}.
\]
For a non-realized pair the weight is \(0\), so the same inequality holds.

The event bound is also pointwise.  If
\(\omega\in T\cap H\cap G_{B,J,\beta,\tau}\), then the non-realized
exclusion makes \((B,J,\beta,\tau)\) realized.  The cell geometry says that
\(\omega\) satisfies all present and absent constraints of the terminal
cylinder defining \(w_{B,J,\beta,\tau}\).  Hence the conditional probability
of \(T\cap G_{B,J,\beta,\tau}\) given \(H\) is at most this cell weight.  To
convert this conditional bound into an unconditional event bound, use
\(\noderef{TwoBiteSampleWeightNonnegative}\), which gives nonnegative sample
weights from \(0\le p\le1/2\), hence \(p\le1\).  If \(\Pr(H)=0\), monotonicity
and nonnegativity give
\(\Pr(T\cap H\cap G_{B,J,\beta,\tau})=0\).  If \(\Pr(H)>0\), the definition
of conditional probability gives the same inequality after multiplying by
\(\Pr(H)\).  This is exactly the zero-history-mass split packaged in
\(\noderef{TwoBiteConditionalIntersectionBound}\), applied with cell bound
\(\Pr(H)\) and conditional bound \(w_{B,J,\beta,\tau}\).  Since
\(w_{B,J,\beta,\tau}\ge0\), the resulting \( \min(1,w_{B,J,\beta,\tau})\)
factor is at most \(w_{B,J,\beta,\tau}\), and therefore
\[
\Pr(T\cap H\cap G_{B,J,\beta,\tau})
\le \Pr(H)\,w_{B,J,\beta,\tau}.
\]

Fix \(B\in{\mathcal B}\) and \(J\).  The zero rule from
\(\noderef{FixedSetHistoryCellRedBranchTranscriptCells}\) removes all
non-realized formal pairs from both the probability and weight sums, so it
remains to sum over realized \((\beta,\tau)\).  This zero rule is used only
for removing empty formal cells; the geometry for the remaining realized
pairs comes from the separate realized-geometry certificate exported by the
same node.

For a realized pair, that realized-geometry certificate from
\(\noderef{FixedSetHistoryCellRedBranchTranscriptCells}\) gives
\[
B\cup R_{\beta,J}\subseteq P_\tau,\qquad |B|=u_B,\qquad
|R_{\beta,J}|=u_R,
\]
with \(B\) blue-only and \(R_{\beta,J}\) red-only, hence disjoint.  It also
gives
\[
Z_\tau\subseteq A_\tau,\qquad P_\tau\cap A_\tau=\emptyset,\qquad
L\le |Z_\tau|.
\]

For this fixed \(B,J\), choose the residual leaf certificate assumed in the
statement.  Its nonnegativity and normalized total mass give a genuine
leaf-level subprobability; its injectivity says that no realized
branch/transcript pair is counted through the same residual leaf twice; and
its termwise domination is exactly
\[
p^{|P_\tau|-u_R-u_B}q^{|A_\tau|-|Z_\tau|}
\le \nu_{B,J}(\iota_{B,J}(\beta,\tau)).
\]
The certificate is explicitly allowed to charge selected blue absences:
\(Z_\tau\) is only required to be a subset of \(A_\tau\), and its factors are
removed once from the same terminal cylinder.

Applying
\(\noderef{FixedSetHistoryCellFixedBJResidualProductTreePacking}\) to this
same residual leaf certificate gives residual masses \(\nu_{\beta,\tau}\)
satisfying
\[
\sum_{\beta,\tau}\nu_{\beta,\tau}\le1
\]
and, for every realized pair,
\[
p^{|P_\tau|}q^{|A_\tau|}
\le
p^{u_R}p^{u_B}q^L\,\nu_{\beta,\tau}.
\]

Combining the helper's termwise inequality with the cylinder estimate already
proved,
\[
w_{B,J,\beta,\tau}\le p^{|P_\tau|}q^{|A_\tau|},
\]
gives
\[
w_{B,J,\beta,\tau}
\le p^{u_R}p^{u_B}q^L\,\nu_{\beta,\tau}
\]
for every realized pair.  Non-realized formal pairs have weight \(0\), so the
same domination holds after setting their residual mass to \(0\).

Summing and using nonnegativity of the fixed factor gives
\[
\sum_{\beta,\tau} w_{B,J,\beta,\tau}
\le
p^{u_R}p^{u_B}q^L
  \sum_{\beta,\tau}\nu_{\beta,\tau}
\le
p^{u_R}p^{u_B}q^L .
\]
Substituting \(q=1-p\) gives the claimed bound.
\end{proof}
